\begin{longtable}[l]{ |m{8.2cm}|r|c|r|m{9.3cm}| }
\caption{BLOCK0 参数}
    \label{tab:efuse-block0-para}
    \\\hline
    \rowcolor{lightgray}
        \textbf{参数} &  \textbf{\thead{位宽}} & \textbf{\thead{硬件使用}} &  \textbf{\thead{\hyperref[fielddesc:EFUSEWRDIS]{EFUSE\_WR\_DIS} \\烧写保护位}}
    & \textbf{描述} \\
    \endfirsthead
    \multicolumn{5}{c}{\bfseries \tablename\ \thetable{} -- 接上页}\\\hline
    \rowcolor{lightgray}
        \textbf{参数} & \textbf{\thead{位宽}} & \textbf{\thead{硬件使用}} &  \textbf{\thead{\hyperref[fielddesc:EFUSEWRDIS]{EFUSE\_WR\_DIS} \\烧写保护位}}
    & \textbf{描述} \\
    \endhead
    \multicolumn{5}{r}{\textbf{见下页}} \\
    \endfoot
    \endlastfoot
    \hline
    \textbf{\hyperref[fielddesc:EFUSEWRDIS]{EFUSE\_WR\_DIS}}  & 32 & Y & N/A & 表示是否禁止 eFuse 控制器烧写对应参数 \\
    \hline
    \textbf{\hyperref[fielddesc:EFUSERDDIS]{EFUSE\_RD\_DIS}}  & 7 & Y & 0  & 表示是否禁止用户读取 eFuse 存储器中 BLOCK4 \verb+~+ 10 的内容 \\
    \hline
    \hyperref[fielddesc:EFUSEUSBDEVICEEXCHGPINS]{EFUSE\_USB\_DEVICE\_EXCHG\_PINS} & 1 & Y & 30 & 表示是否交换 USB Serial/JTAG 的 D+/D- 管脚 \\
    \hline
    \hyperref[fielddesc:EFUSEUSBOTG11EXCHGPINS]{EFUSE\_USB\_OTG11\_EXCHG\_PINS} & 1 & Y & 30 & 表示是否交换 USB FS D+/D- 管脚 \\
    \hline
    \hyperref[fielddesc:EFUSEDISUSBJTAG]{EFUSE\_DIS\_USB\_JTAG}  & 1 & Y & 2  & 表示是否关闭 USB Serial/JTAG 模块中 USB 转 JTAG 的功能 \\
    \hline
    \hyperref[fielddesc:EFUSEPOWERGLITCHEN]{EFUSE\_POWERGLITCH\_EN}  & 1 & Y & 2  & 表示使能电源毛刺检测功能 \\
    \hline
    \hyperref[fielddesc:EFUSEDISFORCEDOWNLOAD]{EFUSE\_DIS\_FORCE\_DOWNLOAD} & 1 & Y & 2 & 表示是否关闭强制芯片进入 Download 模式的功能\\
    \hline
    \hyperref[fielddesc:EFUSESPIDOWNLOADMSPIDIS]{EFUSE\_SPI\_DOWNLOAD\_MSPI\_DIS} & 1 & Y & 2 & 表示在 boot\_mode\_download 过程中是否关闭 SYS AXI 对 MSPI flash/MSPI RAM 的访问\\
    \hline
    \hyperref[fielddesc:EFUSEDISTWAI]{EFUSE\_DIS\_TWAI} & 1 & Y & 2 & 表示是否关闭 TWAI 控制器功能 \\
    \hline
    \hyperref[fielddesc:EFUSEJTAGSELENABLE]{EFUSE\_JTAG\_SEL\_ENABLE} & 1 & Y & 2 & 表示 \hyperref[fielddesc:EFUSEDISPADJTAG]{EFUSE\_DIS\_PAD\_JTAG} 和 \hyperref[fielddesc:EFUSEDISUSBJTAG]{EFUSE\_DIS\_USB\_JTAG} 均配置为 0 时，是否使能通过 strapping 值选择 JTAG 的信号源\\
    \hline
    \hyperref[fielddesc:EFUSESOFTDISJTAG]{EFUSE\_SOFT\_DIS\_JTAG} & 3 & Y & 31 & 表示是否软禁用 PAD JTAG\\
    \hline
    \hyperref[fielddesc:EFUSEDISPADJTAG]{EFUSE\_DIS\_PAD\_JTAG} & 1 & Y & 2 & 表示是否硬禁用 PAD JTAG（永久） \\
    \hline
    \hyperref[fielddesc:EFUSEDISDOWNLOADMANUALENCRYPT]{EFUSE\_DIS\_DOWNLOAD\_MANUAL\_ENCRYPT} &  1 & Y & 2 & 表示是否禁用 flash 加密功能（SPI 启动模式除外）\\
    \hline
    \hyperref[fielddesc:EFUSEUSBPHYSEL]{EFUSE\_USB\_PHY\_SEL} &  1 & Y & 30 & 表示 USB Serial/JTAG 控制器和 OTG\_FS 与 FS\_PHY1 和 FS\_PHY2 的连接关系\\
    \hline
    \hyperref[fielddesc:EFUSEXTSKEYLENGTH256]{EFUSE\_XTS\_KEY\_LENGTH\_256} &  1 & N & 1 & 表示 flash 加密的密钥\\
    \hline
    \hyperref[fielddesc:EFUSEWDTDELAYSEL]{EFUSE\_WDT\_DELAY\_SEL} & 2 & Y & 2 & 选择 RTC 看门狗 STG0 超时阈值的档位 \\
    \hline
    \hyperref[fielddesc:EFUSESPIBOOTCRYPTCNT]{EFUSE\_SPI\_BOOT\_CRYPT\_CNT} & 3 & Y & 4 & 表示是否使能 SPI boot 加解密\\
    \hline
    \hyperref[fielddesc:EFUSESECUREBOOTKEYREVOKE0]{EFUSE\_SECURE\_BOOT\_KEY\_REVOKE0} & 1 & N & 5 & 表示是否撤销第一个安全启动 (Secure Boot) 密钥 \\
    \hline
    \hyperref[fielddesc:EFUSESECUREBOOTKEYREVOKE1]{EFUSE\_SECURE\_BOOT\_KEY\_REVOKE1} & 1 & N & 6 & 表示是否撤销第二个安全启动密钥 \\
    \hline
    \hyperref[fielddesc:EFUSESECUREBOOTKEYREVOKE2]{EFUSE\_SECURE\_BOOT\_KEY\_REVOKE2} & 1 & N & 7 & 表示是否撤销第三个安全启动密钥 \\
    \hline
    \hyperref[fielddesc:EFUSEKEYPURPOSE0]{EFUSE\_KEY\_PURPOSE\_0} & 4 & Y & 8 & {表示 Key0 用途 (purpose)，见表 \ref{tab:efuse-key-purpose}} \\
    \hline
    \hyperref[fielddesc:EFUSEKEYPURPOSE1]{EFUSE\_KEY\_PURPOSE\_1} & 4 & Y & 9 & {表示 Key1 用途，见表 \ref{tab:efuse-key-purpose}} \\
    \hline
    \hyperref[fielddesc:EFUSEKEYPURPOSE2]{EFUSE\_KEY\_PURPOSE\_2} & 4 & Y & 10 & {表示 Key2 用途，见表 \ref{tab:efuse-key-purpose}} \\
    \hline
    \hyperref[fielddesc:EFUSEKEYPURPOSE3]{EFUSE\_KEY\_PURPOSE\_3} & 4 & Y & 11 & {表示 Key3 用途，见表 \ref{tab:efuse-key-purpose}} \\
    \hline
    \hyperref[fielddesc:EFUSEKEYPURPOSE4]{EFUSE\_KEY\_PURPOSE\_4} & 4 & Y & 12 & {表示 Key4 用途，见表 \ref{tab:efuse-key-purpose}} \\
    \hline
    \hyperref[fielddesc:EFUSEKEYPURPOSE5]{EFUSE\_KEY\_PURPOSE\_5} & 4 & Y & 13 & {表示 Key5 用途，见表 \ref{tab:efuse-key-purpose}} \\
    \hline
    \hyperref[fielddesc:EFUSESECDPALEVEL]{EFUSE\_SEC\_DPA\_LEVEL} & 2 & Y & 14 & 表示防差分功耗分析 (differential power analysis, DPA) 攻击的安全级别 \\
    \hline
    \hyperref[fielddesc:EFUSECRYPTDPAENABLE]{EFUSE\_CRYPT\_DPA\_ENABLE} & 1 & Y & 14 & 表示开启防 DPA 攻击功能 \\
    \hline
    \hyperref[fielddesc:EFUSESECUREBOOTEN]{EFUSE\_SECURE\_BOOT\_EN} & 1 & N & 15 & 表示是否使能安全启动 \\
    \hline
    \hyperref[fielddesc:EFUSESECUREBOOTAGGRESSIVEREVOKE]{EFUSE\_SECURE\_BOOT\_AGGRESSIVE\_REVOKE} & 1 & N & 16 & 表示是否使 Secure Boot 的撤销采用激进策略 \\
    \hline
    \hyperref[fielddesc:EFUSEFLASHTYPE]{EFUSE\_FLASH\_TYPE} & 1 & N & 18 & 表示接口 flash 的类型 \\
    \hline
    \hyperref[fielddesc:EFUSEFLASHPAGESIZE]{EFUSE\_FLASH\_PAGE\_SIZE} & 2 & N & 18 & 表示 flash 页大小 \\
    \hline
    \hyperref[fielddesc:EFUSEFLASHECCEN]{EFUSE\_FLASH\_ECC\_EN} & 1 & N & 18 & 表示是否在 flash boot 期间使能 ECC \\
    \hline
    \hyperref[fielddesc:EFUSEDISUSBOTGDOWNLOADMODE]{EFUSE\_DIS\_USB\_OTG\_DOWNLOAD\_MODE} & 1 & N & 18 & 表示是否禁用通过 USB－OTG 的下载模式 \\
    \hline
    \hyperref[fielddesc:EFUSEFLASHTPUW]{EFUSE\_FLASH\_TPUW} & 4 & N & 18 & 表示上电后 flash 等待时间 \\
    \hline
    \hyperref[fielddesc:EFUSEDISDOWNLOADMODE]{EFUSE\_DIS\_DOWNLOAD\_MODE} & 1 & N & 18 & 表示是否关闭所有 Download 模式\\
    \hline
    \hyperref[fielddesc:EFUSEDISDIRECTBOOT]{EFUSE\_DIS\_DIRECT\_BOOT} & 1 & N & 18 & 表示是否关闭 direct boot 模式\\
    \hline
    \hyperref[fielddesc:EFUSEDISUSBSERIALJTAGROMPRINT]{EFUSE\_DIS\_USB\_SERIAL\_JTAG\_ROM\_PRINT} & 1 & N & 18 & 表示是否关闭 ROM boot 过程中的 USB-Serial-JTAG 打印 \\
    \hline
    \hyperref[fielddesc:EFUSEDISUSBSERIALJTAGDOWNLOADMODE]{EFUSE\_DIS\_USB\_SERIAL\_JTAG\_DOWNLOAD\_MODE} & 1 & N & 18 & 表示是否禁用 USB-Serial-JTAG 下载模式 \\
    \hline
    \hyperref[fielddesc:EFUSEENABLESECURITYDOWNLOAD]{EFUSE\_ENABLE\_SECURITY\_DOWNLOAD} & 1 & N & 18 & 表示是否使能安全下载  \\
    \hline
    \hyperref[fielddesc:EFUSEUARTPRINTCONTROL]{EFUSE\_UART\_PRINT\_CONTROL} & 2 & N & 18 & 表示 UART 打印模式\\
    \hline
    \hyperref[fielddesc:EFUSEFORCESENDRESUME]{EFUSE\_FORCE\_SEND\_RESUME} & 1 & N & 18 & 表示是否强制 ROM 代码在 SPI 启动过程中发送恢复指令 \\
    \hline
    \hyperref[fielddesc:EFUSESECUREVERSION]{EFUSE\_SECURE\_VERSION} & 16 & N & 18 & 表示安全版本，用于 ESP-IDF 的防回滚功能\\
    \hline
    \hyperref[fielddesc:EFUSESECUREBOOTDISABLEFASTWAKE]{EFUSE\_SECURE\_BOOT\_DISABLE\_FAST\_WAKE} & 1 & N & 18 & 表示当 Secure Boot 开启时是否关闭 FAST VERIFY ON WAKE\\
    \hline
    \hyperref[fielddesc:EFUSEHYSENPAD]{EFUSE\_HYS\_EN\_PAD} & 1 & Y & 2 & 表示是否对 PAD0 \verb+~+ 27 使能迟滞功能\\
    \hline
    \hyperref[fielddesc:EFUSEDCDCVSET]{EFUSE\_DCDC\_VSET} & 5 & Y & N/A & 表示 DCDC 的默认电压值\\
    \hline
    \hyperref[fielddesc:EFUSE0PXATIEHSEL0]{EFUSE\_0PXA\_TIEH\_SEL\_0} & 2 & Y & 2 & 表示 LDO VO0 的供电源头选择\\
    \hline
    \hyperref[fielddesc:EFUSE0PXATIEHSEL1]{EFUSE\_0PXA\_TIEH\_SEL\_1} & 2 & Y & 2 & 表示 LDO VO1 的供电源头选择\\
    \hline
    \hyperref[fielddesc:EFUSE0PXATIEHSEL2]{EFUSE\_0PXA\_TIEH\_SEL\_2} & 2 & Y & 2 & 表示 LDO VO2 的供电源头选择\\
    \hline
    \hyperref[fielddesc:EFUSE0PXATIEHSEL3]{EFUSE\_0PXA\_TIEH\_SEL\_3} & 2 & Y & 2 & 表示 LDO VO3 的供电源头选择\\
    \hline
    \hyperref[fielddesc:EFUSEHPPWRSRCSEL]{EFUSE\_HP\_PWR\_SRC\_SEL} & 1 & Y & 3 & 表示 HP 系统电源域选择\\
    \hline
    \hyperref[fielddesc:EFUSEDCDCVSETEN]{EFUSE\_DCDC\_VSET\_EN} & 1 & Y & N/A & 表示是否使用 \hyperref[fielddesc:EFUSEDCDCVSET]{EFUSE\_DCDC\_VSET} 配置的默认电压\\
    \hline
    \hyperref[fielddesc:EFUSEDISWDT]{EFUSE\_DIS\_WDT} & 1 & Y & 2 & 表示是否禁用看门狗功能\\
    \hline
    \hyperref[fielddesc:EFUSEDISSWD]{EFUSE\_DIS\_SWD} & 1 & Y & 2 & 表示是否禁用超级看门狗功能\\
    \hline
\end{longtable}
